%%================================================
%% Chapter 3 - System Design and Development
%%================================================
\chapter{วิธีการออกแบบและพัฒนาระบบ}
\label{chapter3}

\section{ภาพรวมของระบบ}
\label{sec:ch3_system_design}
ระบบต้นแบบ WheelSense ถูกออกแบบเป็นระบบดูแลและติดตามสถานะผู้ใช้รถเข็นและผู้สูงอายุในสภาพแวดล้อมภายในอาคาร โดยรวมข้อมูลจากเซ็นเซอร์รถเข็น อุปกรณ์สวมใส่ สมาร์ทโฟนที่ทำหน้าที่เป็น mobile gateway เซิร์ฟเวอร์ส่วนกลาง แดชบอร์ดตามบทบาทผู้ใช้ และระบบผู้ช่วย AI/MCP ให้อยู่ในโครงสร้างเดียวกัน จุดสำคัญของการออกแบบคือการลดภาระของผู้ดูแลที่ต้องตรวจข้อมูลจากหลายแหล่งด้วยตนเอง และทำให้ข้อมูลตำแหน่ง การเคลื่อนไหว สัญญาณชีพ และเหตุการณ์ที่เกี่ยวข้องถูกส่งเข้าสู่ backend เดียวกันก่อนแสดงผลหรือใช้ประกอบการวิเคราะห์

แนวทางของระบบไม่ได้เริ่มจากการนำไปใช้งานจริงในสถานดูแลผู้สูงอายุทันที แต่ปรับรูปแบบจากการทดสอบภายในพื้นที่ทดลองของโครงงาน เมื่อพบข้อจำกัดของการติดตั้งโหนด การเชื่อมต่อ Wi-Fi การรับค่า RSSI และความล่าช้าของการเรียกใช้ AI จึงปรับสถาปัตยกรรมให้สมาร์ทโฟนเป็น gateway หลักแทนการให้อุปกรณ์ภาคสนามส่งข้อมูลขึ้นระบบโดยตรงทั้งหมด สถาปัตยกรรมที่ใช้ในเล่มนี้จึงแยกภาระงานออกเป็น 4 ส่วน คือ field devices, mobile gateway, platform services และ dashboard/AI ดังรูปที่~\ref{fig:ch3-system-architecture}

\begin{figure}[!htbp]
    \centering
    \resizebox{0.97\textwidth}{!}{%
    \begin{tikzpicture}[
        title/.style={align=center, font=\bfseries\small},
        box/.style={draw, rounded corners, align=center, text width=2.95cm, minimum height=0.76cm, font=\scriptsize, fill=white},
        arrow/.style={->, >=Stealth, thick},
        back/.style={->, >=Stealth, thick, dashed}
    ]
        \node[title] at (0,2.95) {Field devices};
        \node[title] at (4.0,2.95) {Mobile gateway};
        \node[title] at (8.0,2.95) {Platform services};
        \node[title] at (12.0,2.95) {Dashboard / AI};

        \node[box, fill=orange!12] (m5) at (0,2.15) {M5StickC Plus2\\BLE wheelchair sensor};
        \node[box, fill=orange!12] (polar) at (0,1.10) {Polar Verity Sense\\HR, RR interval, PPG};
        \node[box, fill=orange!12] (beacon) at (0,0.05) {BLE RSSI nodes\\room references};

        \node[box, fill=green!12, text width=3.05cm, minimum height=1.15cm] (phone) at (4.0,1.10) {Flutter mobile gateway\\BLE pairing, RSSI scan\\step counter, MQTT/REST};

        \node[box, fill=blue!10] (mqtt) at (8.0,2.15) {MQTT broker\\telemetry topics};
        \node[box, fill=blue!10] (api) at (8.0,1.10) {FastAPI backend\\validate, query, auth};
        \node[box, fill=blue!10] (db) at (8.0,0.05) {PostgreSQL 16\\structured + JSONB};

        \node[box, fill=purple!10] (dash) at (12.0,1.68) {Next.js dashboard\\role-based HMI};
        \node[box, fill=purple!10] (ai) at (12.0,0.52) {AI/MCP pipeline\\RAG, tools, gate};

        \draw[arrow] (m5.east) -- node[above, font=\tiny] {BLE GATT} (phone.west |- m5.east);
        \draw[arrow] (polar.east) -- node[above, font=\tiny] {Polar SDK/BLE} (phone.west);
        \draw[arrow] (beacon.east) -- node[below, font=\tiny] {RSSI scan} (phone.west |- beacon.east);

        \draw[arrow] (phone.east) -- node[above, font=\tiny] {MQTT telemetry} (mqtt.west);
        \draw[arrow] (phone.east) -- node[above, font=\tiny] {REST/config} (api.west);
        \draw[arrow] (mqtt.south) -- (api.north);
        \draw[arrow] (api.south) -- (db.north);
        \draw[back] (api.west) -- ++(-0.70,0) |- node[pos=0.24, left, font=\tiny] {config/command} (phone.south);

        \draw[arrow] (api.east) -- node[above, font=\tiny] {REST/WS} (dash.west);
        \draw[arrow] (db.east) -- node[below, font=\tiny] {context} (ai.west);
        \draw[back] (ai.south) |- ++(-1.65,-0.52) -| node[pos=0.72, below, font=\tiny] {tool request} (api.south);
    \end{tikzpicture}}
    \caption{สถาปัตยกรรมโดยรวมของระบบ WheelSense จากอุปกรณ์ภาคสนามไปยัง mobile gateway, platform services, dashboard และ AI}
    \label{fig:ch3-system-architecture}
\end{figure}

ข้อมูลในระบบเริ่มจาก M5StickC Plus2 ที่ติดตั้งบนรถเข็นและทำหน้าที่เป็น BLE sensor สำหรับอ่านข้อมูล IMU และแบตเตอรี่ อุปกรณ์ Polar Verity Sense ส่งข้อมูลสัญญาณชีพผ่าน Polar SDK และสมาร์ทโฟนอ่าน RSSI จาก BLE node สำหรับใช้ประกอบการระบุตำแหน่งภายในอาคาร ข้อมูลเหล่านี้ถูกรวมที่แอป Flutter gateway ก่อนส่งเข้าสู่ MQTT broker หรือ REST API ตามชนิดข้อมูล จากนั้น FastAPI จะตรวจสอบรูปแบบข้อมูล บันทึกลง PostgreSQL และเปิดให้ dashboard หรือ AI pipeline ดึงไปใช้งานต่อ

\section{การเชื่อมต่ออุปกรณ์ภาคสนามและ Mobile Gateway}
\label{sec:ch3_field_gateway}
การออกแบบส่วนเชื่อมต่ออุปกรณ์ภาคสนามผ่านการทดลองหลายรูปแบบเพื่อหาจุดสมดุลระหว่างความง่ายในการติดตั้ง ความต่อเนื่องของข้อมูล และภาระของ backend ในระยะแรกของแนวคิด ผู้พัฒนาเคยใช้โครงสร้าง Wi-Fi mesh โดยให้ Node T-SimCam รับข้อมูลจาก Xiao nRF52840 ที่โฆษณาสัญญาณ BLE แล้วส่งต่อข้อมูลเป็นทอด ๆ ไปยัง gateway node ที่มีอินเทอร์เน็ต ขณะเดียวกันฝั่งเว็บเริ่มจาก Vite ร่วมกับ React เพราะตอบสนองเร็วและพัฒนา UI ได้ง่าย แต่เมื่อระบบต้องรับ telemetry หลายชนิดและมี workflow เพิ่มขึ้น backend ต้องรับภาระจัดรูปแบบข้อมูลจำนวนมากเกินไป ข้อจำกัดที่พบในระยะนี้คือการติดตั้ง mesh ใช้เวลานาน การหา route ในพื้นที่จริงทำได้ยากเมื่อสภาพแวดล้อมเปลี่ยน และ delay สูงเกินกว่าจะใช้ติดตามสถานะผู้ใช้แบบต่อเนื่องได้ดี

แนวทางต่อมาคือย้ายภาระ gateway บางส่วนมาไว้ที่ M5StickC Plus2 โดยให้อุปกรณ์เชื่อมต่อ Wi-Fi และส่งค่า RSSI จาก node รวมถึงข้อมูลเซ็นเซอร์ของรถเข็นขึ้น server เอง เมื่อไม่มี Wi-Fi จึงเพิ่ม Wi-Fi Manager ให้ M5StickC เปิดเป็น access point เพื่อให้โทรศัพท์เข้าไปตั้งค่าและใช้ hotspot ได้ ส่วน node ที่ต้องส่งภาพหรือข้อมูลจากอาคารยังใช้เครือข่ายของพื้นที่และสื่อสารผ่าน MQTT/HTTP เพื่อส่งสถานะ คำสั่งควบคุม และภาพจากกล้องเข้าสู่ server ในช่วงนี้ frontend ถูกย้ายไปใช้ Next.js เพื่อรองรับ routing, server-side data และโครงสร้าง dashboard ที่ซับซ้อนขึ้น วิธีนี้ลดการพึ่งพา mesh ลง แต่ยังมีข้อจำกัดสำคัญคือเมื่อเครือข่าย Wi-Fi ไม่พร้อมใช้งาน การส่งข้อมูลจะหยุดชะงัก

เมื่อเพิ่ม Polar Verity Sense เข้ามาในระบบ ข้อมูล heart rate, RR interval และ PPG ต้องอ่านผ่าน Polar SDK ทำให้โทรศัพท์เริ่มมีบทบาทสำคัญมากขึ้น มีการพัฒนา mobile app ต้นแบบด้วย React Native และ Expo ซึ่งสามารถเชื่อมต่อ sensor และส่งข้อมูลผ่าน MQTT ได้ แต่การ build และ debug บนอุปกรณ์จริงใช้เวลานาน อีกทั้ง M5StickC Plus2 ยังทำหน้าที่ทั้งรับ RSSI จาก node และส่งข้อมูลผ่าน Wi-Fi ทำให้ตำแหน่งของผู้ใช้มีข้อมูลซ้ำจากสองแหล่ง คือ RSSI ที่ M5StickC รับได้และ RSSI ที่โทรศัพท์รับได้ การจัดการข้อมูลจึงเริ่มซับซ้อนและเสี่ยงต่อการตีความตำแหน่งผิด

จากข้อจำกัดดังกล่าว ระบบที่ใช้ในต้นแบบปัจจุบันจึงให้สมาร์ทโฟนเป็น mobile gateway หลักเพียงจุดเดียว โดย M5StickC Plus2 ลดบทบาทเหลือเป็น BLE wheelchair sensor ที่ส่งข้อมูลไปยังโทรศัพท์โดยตรง Polar Verity Sense เชื่อมต่อกับโทรศัพท์ผ่าน Polar SDK และโทรศัพท์ยังอ่านจำนวนก้าวจาก step counter เพื่อรองรับผู้ใช้งานที่ไม่ได้อยู่บนรถเข็นตลอดเวลา การจัดวางนี้ทำให้ข้อมูลภาคสนามถูกรวมและจัดรูปแบบก่อนขึ้น server ลดปัญหาการมีข้อมูล RSSI ซ้ำจากหลายแหล่ง และช่วยให้ผู้ดูแลตรวจสถานะการเชื่อมต่อจากหน้า gateway ได้ง่ายขึ้น

\begin{figure}[!htbp]
    \centering
    \resizebox{0.96\textwidth}{!}{%
    \begin{tikzpicture}[
        box/.style={draw, rounded corners, align=left, text width=4.25cm, minimum height=0.9cm, font=\scriptsize, fill=white},
        issue/.style={draw, rounded corners, align=left, text width=5.0cm, minimum height=0.9cm, font=\scriptsize, fill=red!6},
        arrow/.style={->, >=Stealth, thick}
    ]
        \node[box, fill=orange!10] (a) {\textbf{Prototype A}\\T-SimCam Wi-Fi mesh, Xiao BLE advertise, Vite + React};
        \node[issue, right=0.7cm of a] (ai) {ติดตั้งนาน หา route ยาก delay สูง และ backend รับภาระจัดรูปแบบข้อมูลมาก};
        \node[box, fill=yellow!12, below=0.42cm of a] (b) {\textbf{Prototype B}\\M5StickC Plus2 เป็น Wi-Fi gateway, RSSI + wheelchair sensor, Next.js};
        \node[issue, right=0.7cm of b] (bi) {ไม่มี Wi-Fi แล้วส่งข้อมูลไม่ได้ ต้องพึ่ง Wi-Fi Manager, AP และ hotspot};
        \node[box, fill=yellow!12, below=0.42cm of b] (c) {\textbf{Prototype C}\\Polar Verity Sense, React Native/Expo mobile gateway, MQTT};
        \node[issue, right=0.7cm of c] (ci) {build/debug บนมือถือช้า และเกิด RSSI ซ้ำจาก M5 กับโทรศัพท์};
        \node[box, fill=green!12, below=0.42cm of c] (d) {\textbf{Current Approach}\\Flutter phone gateway หลัก, M5 เป็น BLE wheelchair sensor, Polar + step counter};
        \node[issue, fill=green!6, right=0.7cm of d] (di) {รวมข้อมูลที่โทรศัพท์ก่อนขึ้น server ลด gateway ซ้ำและควบคุมเส้นทางข้อมูลได้ง่ายขึ้น};

        \draw[arrow] (a.south) -- (b.north);
        \draw[arrow] (b.south) -- (c.north);
        \draw[arrow] (c.south) -- (d.north);
        \draw[arrow] (a) -- (ai);
        \draw[arrow] (b) -- (bi);
        \draw[arrow] (c) -- (ci);
        \draw[arrow] (d) -- (di);
    \end{tikzpicture}}
    \caption{ลำดับการปรับแนวทางการเชื่อมต่ออุปกรณ์ภาคสนามจากการทดลองภายในพื้นที่ทดสอบ}
    \label{fig:ch3-system-layers}
\end{figure}
\FloatBarrier

ในต้นแบบ Node T-SimCam ถูกใช้เป็นตัวอย่างของ camera/BLE node ภาคสนาม โค้ดเฟิร์มแวร์ใน \texttt{firmware/Node\_Tsimcam} เริ่มจากโหลด configuration จากหน่วยความจำ ตั้งค่า camera และเริ่ม BLE beacon ด้วยชื่อ node ที่กำหนดไว้ หากยังไม่ได้ตั้งค่าหรือผู้ใช้กดปุ่ม boot ระหว่างเริ่มต้น อุปกรณ์จะเข้าสู่ config portal ผ่าน Wi-Fi access point เพื่อรับค่า device id, node id, Wi-Fi และ MQTT server แต่ถ้ามี configuration แล้วจะเชื่อมต่อ Wi-Fi เปิดหน้า HTTP status/config และเชื่อมต่อ MQTT broker จากนั้น loop หลักจะรับคำสั่งผ่าน MQTT/HTTP ส่ง registration, status, ack และภาพถ่ายหรือ frame ตามคำสั่งที่ได้รับ

\begin{figure}[!htbp]
    \centering
    \resizebox{0.96\textwidth}{!}{%
    \begin{tikzpicture}[
        box/.style={draw, rounded corners, align=center, text width=2.65cm, minimum height=0.78cm, font=\scriptsize, fill=white},
        dec/.style={draw, diamond, aspect=2.25, align=center, inner sep=1.2pt, font=\scriptsize, fill=yellow!12},
        arrow/.style={->, >=Stealth, thick}
    ]
        \node[box, fill=gray!10] (boot) {Boot\\load config};
        \node[box, fill=orange!10, right=0.55cm of boot] (camera) {Init camera\\start BLE beacon};
        \node[dec, right=0.65cm of camera] (setup) {Config\\ready?};
        \node[box, fill=yellow!10, below=0.90cm of setup] (portal) {Config portal\\AP mode + HTTP form};
        \node[box, fill=green!10, right=0.70cm of setup] (wifi) {Connect Wi-Fi\\open HTTP status};
        \node[box, fill=blue!10, right=0.55cm of wifi] (mqtt) {Connect MQTT\\subscribe control};
        \node[box, fill=purple!10, right=0.55cm of mqtt] (loop) {Main loop\\HTTP + MQTT};
        \node[box, fill=purple!10, below=0.90cm of loop] (status) {Registration\\status + ack};
        \node[box, fill=purple!10, left=0.60cm of status] (image) {Capture / stream\\image payload};
        \node[box, fill=red!7, below=0.90cm of mqtt] (retry) {Reconnect\\or return config};

        \draw[arrow] (boot) -- (camera);
        \draw[arrow] (camera) -- (setup);
        \draw[arrow] (setup) -- node[right, font=\tiny] {no} (portal);
        \draw[arrow] (setup) -- node[above, font=\tiny] {yes} (wifi);
        \draw[arrow] (wifi) -- (mqtt);
        \draw[arrow] (mqtt) -- (loop);
        \draw[arrow] (loop) -- (status);
        \draw[arrow] (loop.south west) |- (image.east);
        \draw[arrow] (status.west) -- (image.east);
        \draw[arrow] (image.north) |- (loop.west);
        \draw[arrow] (mqtt.south) -- node[right, font=\tiny] {lost} (retry.north);
        \draw[arrow] (retry.west) -| node[pos=0.25, below, font=\tiny] {retry} (wifi.south);
        \draw[arrow] (portal.west) -- ++(-0.85,0) |- node[pos=0.25, left, font=\tiny] {save/reboot} (boot.south);
    \end{tikzpicture}}
    \caption{วงรอบเฟิร์มแวร์ Node T-SimCam สำหรับ config portal, HTTP status, MQTT status/ack/control และการส่งภาพแบบ capture/stream}
    \label{fig:ch3_node_tsimcam_loop}
\end{figure}
\FloatBarrier

ข้อจำกัดที่พบจากแนวทางนี้คือการเชื่อมต่อ Wi-Fi มีช่วงเวลารอที่ทำให้การเริ่มต้นช้าได้ และการ stream ภาพผ่าน HTTP client ต่อเนื่องอาจกินรอบการทำงานของ node มากกว่าการส่ง telemetry ขนาดเล็ก จึงเหมาะเป็นองค์ประกอบเสริมของต้นแบบมากกว่าการเป็น gateway หลักของระบบสุดท้าย การตัดสินใจย้าย gateway หลักไปไว้ที่โทรศัพท์จึงลดทั้งความซับซ้อนของ route และลดปัญหาความล่าช้าจาก node ที่ต้องทำงานหลายหน้าที่พร้อมกัน

จากลำดับดังกล่าว สถาปัตยกรรมปัจจุบันจึงย้ายจุดรวมข้อมูลมาอยู่ที่โทรศัพท์เป็นหลัก การส่งข้อมูลไม่ให้ไหลจากหลาย gateway พร้อมกันช่วยลดความซ้ำซ้อนของ RSSI และทำให้ backend รับข้อมูลจากแหล่งเดียวที่จัดรูปแบบแล้ว ก่อนจะส่งต่อเข้าสู่ MQTT, REST API, ฐานข้อมูล และ dashboard ดังแสดงในรูปที่~\ref{fig:ch3-gateway-flow}

\begin{figure}[!htbp]
    \centering
    \resizebox{0.95\textwidth}{!}{%
    \begin{tikzpicture}[
        box/.style={draw, rounded corners, align=center, minimum width=2.9cm, minimum height=0.95cm, font=\scriptsize, fill=white},
        arrow/.style={->, >=Stealth, thick}
    ]
        \node[box, fill=orange!10] (s1) {BLE devices\\M5, Polar, beacons};
        \node[box, fill=green!10, right=0.85cm of s1] (s2) {Flutter gateway\\pair, scan, buffer};
        \node[box, fill=blue!10, right=0.85cm of s2] (s3) {MQTT/REST\\publish telemetry};
        \node[box, fill=blue!10, right=0.85cm of s3] (s4) {FastAPI\\validate + store};
        \node[box, fill=purple!10, right=0.85cm of s4] (s5) {Dashboard/AI\\query + action};
        \draw[arrow] (s1) -- node[above, font=\tiny] {BLE} (s2);
        \draw[arrow] (s2) -- node[above, font=\tiny] {JSON payload} (s3);
        \draw[arrow] (s3) -- (s4);
        \draw[arrow] (s4) -- (s5);
    \end{tikzpicture}}
    \caption{ขั้นตอนการเชื่อมต่อรวมระหว่างอุปกรณ์ภาคสนาม mobile gateway และบริการส่วนกลาง}
    \label{fig:ch3-gateway-flow}
\end{figure}
\FloatBarrier

หากมีคำสั่ง config หรือ command ย้อนกลับจากฝั่ง dashboard หรือ AI ระบบจะไม่ส่งคำสั่งลงอุปกรณ์โดยตรงจาก LLM แต่ให้ backend และ mobile gateway ตรวจสอบเงื่อนไขก่อน แล้วจึงเขียนค่าไปยัง BLE characteristic หรือ config ที่เกี่ยวข้อง วิธีนี้ทำให้เส้นทางคำสั่งแยกจากเส้นทาง telemetry หลักและลดความเสี่ยงจากคำสั่งที่เกิดจากบริบทไม่ครบถ้วน

\section{M5StickC Plus2 BLE Wheelchair Sensor}
\label{sec:ch3_wheelchair_sensor}
M5StickC Plus2 ถูกใช้เป็นเซ็นเซอร์ประจำรถเข็นเพื่อตรวจจับการเคลื่อนไหวของตัวรถ ไม่ได้ทำหน้าที่เป็น gateway หลักในสถาปัตยกรรมปัจจุบัน เฟิร์มแวร์ถูกออกแบบให้ทำงานเป็น BLE peripheral พร้อม custom GATT service เพื่อให้แอป Flutter บนโทรศัพท์เชื่อมต่อ อ่านค่า และรับ notification ได้อย่างต่อเนื่อง ข้อมูลที่ส่งออกประกอบด้วยค่า IMU สถานะแบตเตอรี่ metadata ของอุปกรณ์ และค่าที่ประมวลผลแล้ว เช่น ความเร็ว ระยะทาง และความเร่งในช่วงเวลาปัจจุบัน

เมื่อเริ่มทำงาน เฟิร์มแวร์จะโหลด configuration ที่บันทึกไว้ เช่น ชื่ออุปกรณ์ ค่า room หรือค่าจับคู่ จากนั้นเริ่ม BLE advertising เพื่อประกาศ service UUID และรอ mobile gateway เชื่อมต่อ หลังเชื่อมต่อแล้วอุปกรณ์จะอ่านค่า IMU และแบตเตอรี่ตามรอบเวลา ประมวลผล gyro offset เพื่อลด bias ของแกนหมุน ใช้ exponential moving average (EMA) เพื่อลดสัญญาณรบกวน และใช้ deadband เพื่อตัดค่าการเคลื่อนไหวขนาดเล็กที่อาจเกิดจาก noise ก่อนคำนวณความเร็ว ระยะทาง และความเร่งของรถเข็น

\begin{table}[H]
    \centering
    \caption{หน้าที่หลักของเฟิร์มแวร์ M5StickC Plus2 ในระบบปัจจุบัน}
    \label{tab:ch3_m5_hw}
    \small
    \renewcommand{\arraystretch}{1.25}
    \begin{tabular}{|L{4.0cm}|L{8.6cm}|}
        \hline
        \textbf{ส่วนการทำงาน} & \textbf{รายละเอียดการออกแบบ} \\ \hline
        Configuration & โหลดชื่ออุปกรณ์ สถานะจับคู่ ค่า room และค่าที่จำเป็นต่อการเชื่อมต่อ BLE \\ \hline
        BLE GATT service & advertise service และเปิด characteristic สำหรับ telemetry, command, config และ time sync \\ \hline
        Sensor sampling & อ่านค่า acceleration, angular velocity และ battery จากฮาร์ดแวร์ภายใน M5StickC Plus2 \\ \hline
        Signal processing & ชดเชย gyro offset ใช้ EMA และ deadband เพื่อลด noise ก่อนคำนวณ motion \\ \hline
        Telemetry notify & ส่ง JSON payload ผ่าน BLE notification ไปยังโทรศัพท์ที่จับคู่ไว้ \\ \hline
        Command receiver & รับคำสั่ง start, stop, identify, pair, calibrate, reset distance, config และ time sync จากโทรศัพท์ \\ \hline
    \end{tabular}
\end{table}

\begin{figure}[!htbp]
    \centering
    \includegraphics[width=0.72\textwidth]{chapter3/ch3-m5-wheelchair-placement.jpg}
    \caption{ตัวอย่างตำแหน่งติดตั้ง M5StickC Plus2 บนรถเข็นสำหรับอ่านค่าการเคลื่อนไหวและส่งข้อมูลผ่าน BLE ไปยังโทรศัพท์}
    \label{fig:ch3-m5-placement}
\end{figure}

โครงสร้างข้อมูลที่ส่งไปยังโทรศัพท์ถูกกำหนดให้เป็น JSON เพื่อให้ mobile gateway สามารถ normalize และส่งต่อเข้าระบบได้โดยไม่ต้องผูกกับรูปแบบ binary ของเฟิร์มแวร์โดยตรง ตัวอย่างกลุ่มข้อมูลหลัก ได้แก่ timestamp, device identifier, ค่า IMU, motion summary, battery level และสถานะการเชื่อมต่อ การกำหนดข้อมูลเป็น schema ลักษณะนี้ช่วยให้ backend ตรวจสอบและจัดเก็บข้อมูลใน PostgreSQL/JSONB ได้ง่ายขึ้น

\begin{figure}[!htbp]
    \centering
    \resizebox{0.85\textwidth}{!}{%
    \begin{tikzpicture}[
        node distance=0.72cm,
        proc/.style={draw, rounded corners, align=center, minimum width=4.0cm, minimum height=0.78cm, font=\scriptsize, fill=blue!6},
        dec/.style={draw, diamond, aspect=2.1, align=center, inner sep=1pt, font=\scriptsize, fill=yellow!10},
        arrow/.style={->, >=Stealth, thick}
    ]
        \node[proc] (boot) {Boot and load config};
        \node[proc, below=of boot] (adv) {Advertise BLE GATT service};
        \node[dec, below=of adv] (conn) {Phone\\connected?};
        \node[proc, below=of conn] (read) {Read IMU and battery};
        \node[proc, below=of read] (filter) {Gyro offset, EMA, deadband};
        \node[proc, below=of filter] (calc) {Calculate velocity, distance, acceleration};
        \node[proc, below=of calc] (json) {Build JSON telemetry payload};
        \node[proc, below=of json] (notify) {Notify payload to phone};
        \node[dec, below=of notify] (cmd) {Command or\\config write?};
        \node[proc, right=1.25cm of cmd] (apply) {Apply command\\config/time sync};

        \draw[arrow] (boot) -- (adv);
        \draw[arrow] (adv) -- (conn);
        \draw[arrow] (conn) -- node[right, font=\tiny] {yes} (read);
        \draw[arrow] (conn.west) -- ++(-1.2,0) |- node[pos=0.25, left, font=\tiny] {no} (adv.west);
        \draw[arrow] (read) -- (filter);
        \draw[arrow] (filter) -- (calc);
        \draw[arrow] (calc) -- (json);
        \draw[arrow] (json) -- (notify);
        \draw[arrow] (notify) -- (cmd);
        \draw[arrow] (cmd.east) -- node[above, font=\tiny] {yes} (apply.west);
        \draw[arrow] (apply.north) |- (read.east);
        \draw[arrow] (cmd.west) -- ++(-1.2,0) |- node[pos=0.25, left, font=\tiny] {no} (read.west);
    \end{tikzpicture}}
    \caption{แผนผังวงรอบการทำงานของเฟิร์มแวร์ M5StickC Plus2 แบบ BLE wheelchair sensor}
    \label{fig:ch3-firmware-loop}
\end{figure}
\FloatBarrier

\section{Polar Verity Sense และ Flutter Mobile Gateway}
\label{sec:ch3_mobile_gateway}
อุปกรณ์ Polar Verity Sense ถูกเพิ่มเข้ามาเพื่อให้ระบบมีข้อมูลสัญญาณชีพประกอบกับข้อมูลการเคลื่อนไหวของรถเข็น การอ่านข้อมูลจาก Polar ไม่ได้ใช้เพียง BLE GATT มาตรฐานทั่วไปเท่านั้น แต่ใช้ Polar SDK เพื่ออ่านข้อมูลที่ละเอียดขึ้น เช่น heart rate, RR interval และ PPG stream จากนั้นแอป Flutter จะรวมข้อมูลของ Polar เข้ากับข้อมูลจาก M5StickC Plus2 และข้อมูล step counter ของโทรศัพท์

แอป Flutter ทำหน้าที่เป็น gateway หลักในระบบ เพราะเป็นจุดที่มีทั้ง BLE สำหรับเชื่อมต่ออุปกรณ์ภาคสนามและอินเทอร์เน็ตผ่าน Wi-Fi หรือเครือข่ายเซลลูลาร์ หน้าการทำงานที่เกี่ยวข้องกับการติดตั้งภาคสนามคือหน้า \texttt{/WheelSensor} ซึ่งใช้ตรวจสถานะการเชื่อมต่อ BLE, Polar, M5StickC Plus2, MQTT และการส่งข้อมูลขึ้น server จุดนี้ช่วยให้ผู้ดูแลหรือผู้พัฒนาสามารถแยกปัญหาได้ว่าข้อมูลหายจากระดับ pairing, sensor stream, network publish หรือ backend ingest

\begin{figure}[H]
    \centering
    \includegraphics[width=0.82\textwidth]{chapter3/ch3-polar-smartphone.png}
    \caption{การใช้งาน Polar Verity Sense ร่วมกับสมาร์ทโฟนที่ทำหน้าที่เป็น mobile gateway}
    \label{fig:ch3-polar-smartphone}
\end{figure}

\begin{table}[H]
    \centering
    \caption{กลุ่มข้อมูลที่ mobile gateway รวมก่อนส่งเข้าสู่ระบบ}
    \label{tab:ch3_mobile_payload}
    \small
    \renewcommand{\arraystretch}{1.25}
    \begin{tabular}{|L{3.2cm}|L{4.7cm}|L{4.7cm}|}
        \hline
        \textbf{แหล่งข้อมูล} & \textbf{ข้อมูลที่รับเข้า gateway} & \textbf{การส่งต่อเข้าสู่ระบบ} \\ \hline
        M5StickC Plus2 & IMU, motion summary, battery, device metadata & รวมเป็น telemetry ของรถเข็นและ publish ผ่าน MQTT/REST \\ \hline
        Polar Verity Sense & HR, RR interval, PPG และสถานะแบตเตอรี่ตามที่ SDK รองรับ & normalize เป็นข้อมูลสัญญาณชีพของผู้ใช้และส่งพร้อม timestamp \\ \hline
        Phone sensor & step counter และสถานะเครือข่ายของเครื่อง & ใช้ประกอบบริบทผู้ใช้ที่ไม่ได้อยู่บนรถเข็นตลอดเวลา \\ \hline
        BLE RSSI nodes & รายการ beacon ที่พบและค่า RSSI & ส่งเป็นข้อมูลตำแหน่งให้ backend ประเมินห้องหรือโซน \\ \hline
    \end{tabular}
\end{table}

\begin{figure}[!htbp]
    \centering
    \resizebox{0.88\textwidth}{!}{%
    \begin{tikzpicture}[
        node distance=0.75cm,
        proc/.style={draw, rounded corners, align=center, minimum width=4.2cm, minimum height=0.78cm, font=\scriptsize, fill=green!7},
        dec/.style={draw, diamond, aspect=2.2, align=center, inner sep=1pt, font=\scriptsize, fill=yellow!10},
        arrow/.style={->, >=Stealth, thick}
    ]
        \node[proc] (start) {Open gateway service};
        \node[proc, below=of start] (pair) {Scan and pair M5/Polar/beacons};
        \node[proc, below=of pair] (read) {Read BLE streams and step counter};
        \node[proc, below=of read] (pack) {Normalize payload with timestamp};
        \node[dec, below=of pack] (net) {Network\\available?};
        \node[proc, below=of net] (pub) {Publish MQTT or REST};
        \node[proc, below=of pub] (status) {Update WheelSensor status};
        \node[proc, right=1.35cm of net] (buffer) {Buffer and retry\\on reconnect};
        \node[dec, below=of status] (cmd) {Config or\\command?};
        \node[proc, right=1.35cm of cmd] (sendcmd) {Write BLE command\\to M5 if needed};

        \draw[arrow] (start) -- (pair);
        \draw[arrow] (pair) -- (read);
        \draw[arrow] (read) -- (pack);
        \draw[arrow] (pack) -- (net);
        \draw[arrow] (net) -- node[right, font=\tiny] {yes} (pub);
        \draw[arrow] (net.east) -- node[above, font=\tiny] {no} (buffer.west);
        \draw[arrow] (buffer.north) |- (read.east);
        \draw[arrow] (pub) -- (status);
        \draw[arrow] (status) -- (cmd);
        \draw[arrow] (cmd.east) -- node[above, font=\tiny] {yes} (sendcmd.west);
        \draw[arrow] (sendcmd.north) |- (read.east);
        \draw[arrow] (cmd.west) -- ++(-1.1,0) |- node[pos=0.25, left, font=\tiny] {no} (read.west);
    \end{tikzpicture}}
    \caption{วงรอบการทำงานของ Flutter mobile gateway สำหรับรวมข้อมูล Polar, M5StickC Plus2, RSSI และ step counter}
    \label{fig:ch3-smartphone-loop}
\end{figure}

\begin{figure}[!htbp]
    \centering
    \includegraphics[width=0.38\textwidth]{chapter3/ch3-wheelsensor-real.png}
    \caption{หน้าจอ \texttt{/WheelSensor} สำหรับตรวจสถานะ BLE, Polar, M5StickC Plus2 และการส่งข้อมูลขึ้น server}
    \label{fig:ch3_wheelsensor_status}
\end{figure}

\section{Indoor Localization}
\label{sec:ch3_localization}
การระบุตำแหน่งภายในอาคารของระบบใช้ค่า RSSI จาก BLE node ที่ติดตั้งตามพื้นที่อ้างอิง โดย mobile gateway สแกนสัญญาณจาก node รอบตัวและส่งรายการค่า RSSI พร้อม timestamp ไปยัง server สำหรับประเมินห้องหรือโซนของผู้ใช้ การออกแบบในบทนี้ไม่ลงรายละเอียดสมการซ้ำจากบทที่ 2 แต่เน้นกระบวนการนำข้อมูล RSSI เข้าระบบและการเลือกกลยุทธ์ที่ backend รองรับ

ในเชิง implementation ระบบรองรับทั้งการเทียบกับ fingerprint ด้วย KNN และการใช้ max-RSSI เป็น baseline หรือ fallback เมื่อข้อมูล fingerprint ยังไม่พร้อมหรือความมั่นใจต่ำ การใช้ KNN ต้องมีชุดข้อมูลอ้างอิงของแต่ละตำแหน่ง เช่น room, beacon alias, ค่า RSSI และข้อมูลผู้ใช้หรืออุปกรณ์ที่จับคู่ไว้ ส่วน max-RSSI เลือกตำแหน่งจาก beacon ที่มีสัญญาณแรงที่สุดในรอบปัจจุบัน จึงใช้งานง่ายกว่าแต่ไวต่อการสะท้อนและการเปลี่ยนแปลงของสภาพแวดล้อม

\begin{figure}[H]
    \centering
    \resizebox{0.68\textwidth}{!}{%
    \begin{tikzpicture}[
        bleanchor/.style={circle, draw, thick, fill=blue!25, minimum size=0.72cm, font=\bfseries\small},
        zonelabel/.style={font=\bfseries\small},
        note/.style={font=\scriptsize, text=gray}
    ]
        \draw[thick, fill=gray!3] (-2,-1.6) rectangle (6,4.8);
        \draw[thick] (2,-1.6) -- (2,4.8);
        \draw[thick] (-2,1.6) -- (6,1.6);

        \node[zonelabel] at (0,4.35) {Room / Zone A};
        \node[zonelabel] at (4,4.35) {Room / Zone B};
        \node[zonelabel] at (0,1.15) {Room / Zone C};
        \node[zonelabel] at (4,1.15) {Room / Zone D};

        \node[bleanchor] at (0,3.2) {1};
        \node[bleanchor] at (4.0,3.2) {2};
        \node[bleanchor] at (0,0) {3};
        \node[bleanchor] at (4.0,0) {4};

        \node[note] at (2,-2.05) {เส้นแบ่งแทนฉากกั้นหรือขอบเขตพื้นที่ทดลอง};
        \node[note] at (2,-2.45) {หมายเลข 1--4 คือ BLE node ที่วางบริเวณกึ่งกลางของแต่ละ zone};
    \end{tikzpicture}}
    \caption{แผนผังพื้นที่ทดลอง 4 zone พร้อมตำแหน่ง BLE node กึ่งกลางแต่ละพื้นที่สำหรับเก็บค่า RSSI}
    \label{fig:ch3-beacon-layout}
\end{figure}

\begin{figure}[!htbp]
    \centering
    \resizebox{0.95\textwidth}{!}{%
    \begin{tikzpicture}[
        box/.style={draw, rounded corners, align=center, minimum width=2.8cm, minimum height=0.85cm, font=\scriptsize, fill=white},
        arrow/.style={->, >=Stealth, thick}
    ]
        \node[box, fill=orange!10] (scan) {Phone scans\\BLE RSSI nodes};
        \node[box, fill=green!10, right=0.7cm of scan] (send) {Gateway sends\\RSSI list};
        \node[box, fill=blue!10, right=0.7cm of send] (ready) {Backend checks\\assignment + room};
        \node[box, fill=blue!10, right=0.7cm of ready] (knn) {KNN if fingerprint\\is available};
        \node[box, fill=yellow!12, below=0.55cm of knn] (fallback) {Max-RSSI fallback\\or baseline};
        \node[box, fill=purple!10, right=0.7cm of knn] (store) {Store prediction\\and confidence};
        \node[box, fill=purple!10, right=0.7cm of store] (view) {Dashboard/API\\uses location};

        \draw[arrow] (scan) -- (send);
        \draw[arrow] (send) -- (ready);
        \draw[arrow] (ready) -- (knn);
        \draw[arrow] (knn) -- node[above, font=\tiny] {confident} (store);
        \draw[arrow] (knn) -- node[right, font=\tiny] {not ready} (fallback);
        \draw[arrow] (fallback) -| (store);
        \draw[arrow] (store) -- (view);
    \end{tikzpicture}}
    \caption{Flowchart การส่งข้อมูล RSSI และการประเมินตำแหน่งด้วย KNN หรือ max-RSSI fallback}
    \label{fig:ch3-localization-flow}
\end{figure}

การเขียนข้อมูลตำแหน่งกลับเข้าสู่ระบบทำเป็นข้อมูลที่ตรวจสอบย้อนกลับได้ โดยเก็บทั้ง raw RSSI reading และผล room prediction แยกกัน เพื่อให้สามารถวิเคราะห์ภายหลังได้ว่าความคลาดเคลื่อนเกิดจากสัญญาณต้นทาง การตั้งค่า fingerprint หรือการเลือกกลยุทธ์ของ localization ในช่วงเวลานั้น การแยก raw data กับ prediction ยังช่วยให้ Chapter 4 สามารถประเมินความแม่นยำของ KNN และ fallback ได้อย่างเป็นระบบ

\section{Server, API และ Database}
\label{sec:ch3_server}
ส่วน server ของ WheelSense ประกอบด้วย MQTT broker, FastAPI backend, PostgreSQL 16 และบริการที่เกี่ยวข้องกับ dashboard/AI โดยออกแบบให้รองรับข้อมูลทั้งแบบ structured และ semi-structured ข้อมูล structured เช่น users, patients, devices, workspaces, rooms และ assignments ใช้สำหรับกำหนดความสัมพันธ์หลักของระบบ ส่วนข้อมูล semi-structured เช่น telemetry payload, AI trace, device configuration และ event metadata เก็บใน JSONB เพื่อรองรับ payload จากอุปกรณ์ที่อาจเปลี่ยนรายละเอียดระหว่างการทดลอง

MQTT broker รับ telemetry จาก mobile gateway และ topic ที่เกี่ยวกับอุปกรณ์ภาคสนาม ส่วน FastAPI รับทั้งข้อมูลจาก gateway, dashboard และ AI/MCP pipeline ผ่าน REST endpoint หรือ internal service การออกแบบ backend ใช้ async I/O เพื่อให้รับข้อมูลจากหลายอุปกรณ์พร้อมกันได้โดยไม่บล็อกการทำงานของ API หลัก เมื่อข้อมูลเข้ามา server จะทำขั้นตอน ingest, validate, store, query และส่งต่อไปยัง dashboard หรือ AI ตามลำดับ

\begin{table}[H]
    \centering
    \caption{บริการหลักของ platform services ในระบบ WheelSense}
    \label{tab:ch3_services}
    \small
    \renewcommand{\arraystretch}{1.2}
    \begin{tabular}{|L{3.0cm}|L{3.2cm}|L{6.2cm}|}
        \hline
        \textbf{บริการ} & \textbf{เทคโนโลยี} & \textbf{หน้าที่ในระบบ} \\ \hline
        MQTT broker & Eclipse Mosquitto & รับ telemetry จาก mobile gateway และใช้ส่ง config หรือสถานะบางส่วนกลับไปยังอุปกรณ์ \\ \hline
        Backend API & FastAPI & ตรวจสอบข้อมูล จัดการสิทธิ์ ประมวลผล request จาก dashboard/gateway และเปิด API ให้ AI ใช้ข้อมูล \\ \hline
        Database & PostgreSQL 16 + JSONB & เก็บ schema หลักของระบบ telemetry, RSSI, room prediction, AI trace และข้อมูลตั้งค่า \\ \hline
        Frontend & Next.js & แสดงผล dashboard ตามบทบาทผู้ใช้และเรียกข้อมูลจาก backend ผ่าน API \\ \hline
        AI/MCP runtime & Agent runtime + MCP server & จัดเตรียมบริบท ใช้ LLM แบบมีข้อจำกัด และควบคุมการเรียก tool ผ่าน execution gate \\ \hline
    \end{tabular}
\end{table}

\begin{figure}[!htbp]
    \centering
    \resizebox{0.95\textwidth}{!}{%
    \begin{tikzpicture}[
        box/.style={draw, rounded corners, align=center, minimum width=2.8cm, minimum height=0.85cm, font=\scriptsize, fill=white},
        arrow/.style={->, >=Stealth, thick}
    ]
        \node[box, fill=orange!10] (ingest) {Ingest\\MQTT/REST};
        \node[box, fill=yellow!12, right=0.75cm of ingest] (validate) {Validate\\schema + auth};
        \node[box, fill=blue!10, right=0.75cm of validate] (store) {Store\\PostgreSQL/JSONB};
        \node[box, fill=green!10, right=0.75cm of store] (query) {Query\\API/service layer};
        \node[box, fill=purple!10, right=0.75cm of query] (consume) {Dashboard/AI\\view + action};
        \draw[arrow] (ingest) -- (validate);
        \draw[arrow] (validate) -- (store);
        \draw[arrow] (store) -- (query);
        \draw[arrow] (query) -- (consume);
    \end{tikzpicture}}
    \caption{ลำดับการประมวลผลข้อมูลฝั่ง server ตั้งแต่ ingest, validate, store, query ไปจนถึง dashboard และ AI}
    \label{fig:ch3-server-ingestion}
\end{figure}
\FloatBarrier

schema หลักของระบบสามารถแบ่งได้เป็น 5 กลุ่ม ได้แก่ กลุ่ม workspace และสิทธิ์ผู้ใช้ กลุ่ม patient และ assignment กลุ่ม device และ room binding กลุ่ม telemetry/localization และกลุ่ม AI/runtime log ตารางเช่น \texttt{users}, \texttt{patients}, \texttt{devices}, \texttt{workspaces}, \texttt{rooms}, \texttt{imu\_telemetry}, \texttt{rssi\_readings}, \texttt{room\_predictions}, \texttt{mobile\_device\_telemetry}, \texttt{pipeline\_events} และ \texttt{behavioral\_states} ช่วยให้ข้อมูลจากหลายส่วนถูกอ้างอิงกลับไปยังผู้ใช้ อุปกรณ์ และพื้นที่ได้อย่างเป็นระบบ

\begin{figure}[!htbp]
    \centering
    \includegraphics[width=0.98\textwidth]{chapter3/ER-Diagram.png}
    \caption{แผนภาพ schema ฐานข้อมูลที่ใช้เชื่อม workspaces, users, patients, devices และข้อมูล telemetry}
    \label{fig:ch3-database-erd}
\end{figure}
\FloatBarrier

\begin{table}[H]
    \centering
    \caption{ตัวอย่างกลุ่มข้อมูลที่จัดเก็บใน PostgreSQL และ JSONB}
    \label{tab:ch3_database_groups}
    \small
    \renewcommand{\arraystretch}{1.2}
    \begin{tabular}{|L{3.0cm}|L{4.4cm}|L{5.1cm}|}
        \hline
        \textbf{กลุ่มข้อมูล} & \textbf{ตารางหรือ entity ตัวอย่าง} & \textbf{เหตุผลในการออกแบบ} \\ \hline
        User/workspace & users, workspaces, roles & แยกขอบเขตองค์กรและสิทธิ์การเข้าถึงข้อมูลของผู้ใช้แต่ละบทบาท \\ \hline
        Patient/device & patients, devices, assignments & เชื่อมผู้พักอาศัยกับอุปกรณ์ รถเข็น และห้องที่เกี่ยวข้อง \\ \hline
        Telemetry & imu\_telemetry, mobile\_device\_telemetry & เก็บข้อมูลจาก M5, Polar, phone และ gateway พร้อม timestamp \\ \hline
        Localization & rssi\_readings, room\_predictions, training data & แยก raw RSSI ออกจากผลทำนายเพื่อประเมินย้อนหลังได้ \\ \hline
        AI runtime & pipeline\_events, behavioral\_states & เก็บ trace ของแต่ละ layer และ behavioral context สำหรับตรวจสอบการทำงานของ AI \\ \hline
    \end{tabular}
\end{table}

\section{Dashboard และ HMI ตามบทบาทผู้ใช้}
\label{sec:ch3_hmi}
Dashboard ของ WheelSense พัฒนาด้วย Next.js เพื่อเป็นหน้าจอหลักสำหรับผู้ดูแลและผู้ใช้งานตามบทบาท การออกแบบ HMI ไม่ได้ให้ผู้ใช้ทุกกลุ่มเห็นข้อมูลเท่ากัน แต่กำหนด permission และมุมมองตามหน้าที่จริงในระบบ เพื่อให้ผู้ใช้แต่ละกลุ่มเห็นเฉพาะข้อมูลที่จำเป็นต่อการตัดสินใจหรือติดตามสถานะ ขณะที่ mobile app ทำหน้าที่เป็น gateway ภาคสนามสำหรับ pairing sensor และส่งข้อมูล ไม่ใช่ dashboard หลักของทุกบทบาท

\begin{table}[H]
    \centering
    \caption{บทบาทผู้ใช้ หน้าจอหลัก และหน้าที่ที่รองรับบน dashboard}
    \label{tab:ch3_hmi_role_flow}
    \small
    \renewcommand{\arraystretch}{1.2}
    \begin{tabular}{|L{2.4cm}|L{4.4cm}|L{5.8cm}|}
        \hline
        \textbf{บทบาท} & \textbf{หน้าจอหลักที่ใช้} & \textbf{งานที่ทำได้ในระบบ} \\ \hline
        Admin & Dashboard overview, task operations, facility, floorplans, devices, personnel, alerts, messages, support, settings, EaseAI & จัดการ workspace, ผู้ใช้, อุปกรณ์, ห้อง, floorplan, task, profile ของผู้พักอาศัยและ staff รวมถึงตรวจสถานะของ platform และเรียกใช้ AI assistant ภายใต้สิทธิ์ผู้ดูแลระบบ \\ \hline
        Head Nurse & Dashboard overview, alerts, patient profile, task operations, devices, floorplans, personnel, staff, messages, support & ติดตามผู้พักอาศัยหลายคน ตรวจ alert จัดลำดับความสำคัญ มอบหมายงาน และดูข้อมูลสัญญาณชีพ/ตำแหน่งในระดับที่เกี่ยวข้องกับการดูแล \\ \hline
        Supervisor & Dashboard overview, emergency, task operations, floorplans, personnel, prescriptions, patient profile, messages, support & ตรวจภาพรวมงานและเหตุการณ์ในพื้นที่ ติดตาม emergency/task ที่ต้องส่งต่อ และตรวจข้อมูลผู้พักอาศัยหรือบุคลากรในขอบเขตการกำกับดูแล \\ \hline
        Observer & Dashboard overview, alerts, floorplans, prescriptions, task operations, patient profile, personnel, messages, support & เฝ้าดูสถานะและเหตุการณ์แบบจำกัดสิทธิ์ ใช้ข้อมูลเพื่อสังเกตการณ์หรือรายงานต่อผู้ดูแล โดยไม่แก้ไขข้อมูลสำคัญของระบบ \\ \hline
        Patient & Personal dashboard, schedule, pharmacy, room controls, services, messages, support, settings & ดูข้อมูลส่วนตัวของตนเอง เช่น care plan, schedule, medication/service request, room control และช่องทางขอความช่วยเหลือหรือสื่อสารกับผู้ดูแล \\ \hline
    \end{tabular}
\end{table}

ตารางบทบาทข้างต้นถูกแปลงเป็น role flow เพื่อให้เห็นว่าแต่ละกลุ่มผู้ใช้ไม่ได้เข้าถึงข้อมูลชุดเดียวกันทั้งหมด แต่ทุกคำขอจะผ่าน permission และ visibility layer ก่อนเรียกข้อมูลจาก FastAPI หรือ telemetry/AI context รูปที่~\ref{fig:ch3_hmi_role_flow} แสดงลำดับการไหลของหน้าจอจากบทบาทผู้ใช้ไปยัง service หลักของระบบ โดยไม่รวม mobile gateway เพราะ mobile gateway อยู่ในเส้นทางรับส่งข้อมูล sensor ที่อธิบายไว้ก่อนหน้า

\begin{figure}[!htbp]
    \centering
    \resizebox{0.94\textwidth}{!}{%
    \begin{tikzpicture}[
        role/.style={draw, rounded corners, align=center, text width=2.05cm, minimum height=0.56cm, font=\scriptsize, fill=orange!10},
        page/.style={draw, rounded corners, align=center, text width=3.15cm, minimum height=0.6cm, font=\scriptsize, fill=green!10},
        gate/.style={draw, rounded corners, align=center, text width=2.45cm, minimum height=4.7cm, font=\scriptsize, fill=blue!8},
        service/.style={draw, rounded corners, align=center, text width=2.8cm, minimum height=0.65cm, font=\scriptsize, fill=blue!10},
        arrow/.style={->, >=Stealth, thick}
    ]
        \node[role] (admin) at (0,2.4) {Admin};
        \node[role] (hn) at (0,1.2) {Head Nurse};
        \node[role] (sup) at (0,0.0) {Supervisor};
        \node[role] (obs) at (0,-1.2) {Observer};
        \node[role] (pat) at (0,-2.4) {Patient};

        \node[page] (adminp) at (3.55,2.4) {System, task, profile,\\device management};
        \node[page] (hnp) at (3.55,1.2) {Care dashboard,\\alerts, assignment};
        \node[page] (supp) at (3.55,0.0) {Operations and\\emergency overview};
        \node[page] (obsp) at (3.55,-1.2) {Read-only monitoring\\and reporting};
        \node[page] (patp) at (3.55,-2.4) {Own schedule, service,\\room controls};

        \node[gate] (perm) at (7.2,0) {Role permission\\workspace scope\\patient visibility\\action policy};
        \node[service] (api) at (10.8,1.45) {FastAPI\\application API};
        \node[service] (db) at (10.8,0.35) {PostgreSQL/JSONB\\query layer};
        \node[service] (tele) at (10.8,-0.75) {Telemetry, alert\\and location data};
        \node[service] (ai) at (10.8,-1.85) {EaseAI/MCP\\context and tools};

        \draw[arrow] (admin) -- (adminp);
        \draw[arrow] (hn) -- (hnp);
        \draw[arrow] (sup) -- (supp);
        \draw[arrow] (obs) -- (obsp);
        \draw[arrow] (pat) -- (patp);
        \draw[arrow] (adminp.east) -- (perm.west |- adminp.east);
        \draw[arrow] (hnp.east) -- (perm.west |- hnp.east);
        \draw[arrow] (supp.east) -- (perm.west |- supp.east);
        \draw[arrow] (obsp.east) -- (perm.west |- obsp.east);
        \draw[arrow] (patp.east) -- (perm.west |- patp.east);
        \draw[arrow] (perm.east |- api.west) -- (api.west);
        \draw[arrow] (perm.east |- db.west) -- (db.west);
        \draw[arrow] (perm.east |- tele.west) -- (tele.west);
        \draw[arrow] (perm.east |- ai.west) -- (ai.west);
    \end{tikzpicture}}
    \caption{HMI role flow ของระบบ WheelSense แสดงการผ่านสิทธิ์ บริบท workspace และ service ที่ dashboard เรียกใช้งาน}
    \label{fig:ch3_hmi_role_flow}
\end{figure}
\FloatBarrier

การใช้งานที่ซ้ำกันหลายบทบาท เช่น task operations, patient profile, staff profile, messages และ support ถูกอธิบายจากสิทธิ์ของ Admin ก่อน เนื่องจาก Admin เห็นภาพรวมมากที่สุดและใช้ตรวจความครบถ้วนของข้อมูลระบบได้ หน้าจอ task operations ใช้ติดตามงานหรือเหตุการณ์ที่ต้องมอบหมายให้ผู้ดูแล ส่วน patient/staff profile ใช้ดูรายละเอียดรายบุคคลและสถานะที่เกี่ยวข้อง เช่น ข้อมูลผู้พักอาศัย อุปกรณ์ที่ผูกไว้ การแจ้งเตือน และข้อมูลบุคลากรใน workspace เมื่อบทบาทอื่นใช้งานหน้าจอกลุ่มเดียวกัน ระบบจะลดข้อมูลหรือ action ตาม permission ไม่ใช่แยก logic ใหม่ทั้งหมด

\begin{figure}[!htbp]
    \centering
    \includegraphics[width=0.98\textwidth]{chapter3/dashboard-hmi/role-admin-contact.png}
    \caption{ภาพหน้าจอระบบในมุมมอง Admin ครอบคลุม overview, task, profile, device, facility, alert, message, support, setting และ EaseAI}
    \label{fig:ch3_hmi_admin_contact}
\end{figure}
\FloatBarrier

ในมุมมอง Head Nurse ระบบเน้นข้อมูลที่ต้องใช้ดูแลผู้พักอาศัยโดยตรง เช่น alert, patient profile, task operations, device status และ floorplan เพราะเหตุการณ์หนึ่งอาจมาจากอาการหรือพฤติกรรมของผู้พักอาศัย หรืออาจมาจากอุปกรณ์ขาดการเชื่อมต่อ ภาพที่~\ref{fig:ch3_hmi_head_nurse_contact} จึงแสดง workflow หลักที่ Head Nurse ใช้ตรวจสถานะก่อนตัดสินใจส่งงานหรือแจ้งผู้ดูแลคนอื่น

\begin{figure}[!htbp]
    \centering
    \includegraphics[width=0.98\textwidth]{chapter3/dashboard-hmi/role-head-nurse-contact.png}
    \caption{ภาพหน้าจอระบบในมุมมอง Head Nurse สำหรับติดตามผู้พักอาศัย alert, task, device, floorplan และช่องทางสื่อสาร}
    \label{fig:ch3_hmi_head_nurse_contact}
\end{figure}
\FloatBarrier

บทบาท Supervisor และ Observer มีหน้าจอหลายส่วนร่วมกับผู้ดูแล แต่ขอบเขตต่างกัน Supervisor ใช้ dashboard เพื่อมองภาพรวมงาน เหตุการณ์ฉุกเฉิน และข้อมูลที่ต้องส่งต่อในระดับปฏิบัติการ ส่วน Observer ใช้สำหรับเฝ้าดูสถานะหรือบันทึกเหตุการณ์โดยลด action ที่มีผลต่อข้อมูลหลักของระบบ ภาพที่~\ref{fig:ch3_hmi_supervisor_observer_contact} แสดงหน้าจอที่ใช้ในสองบทบาทนี้ โดยคงโครงสร้าง navigation ใกล้กันเพื่อให้ผู้ใช้เปลี่ยนบทบาทหรือฝึกใช้งานได้ง่ายขึ้น

\begin{figure}[!htbp]
    \centering
    \begin{subfigure}[b]{0.49\textwidth}
        \centering
        \includegraphics[width=\linewidth]{chapter3/dashboard-hmi/role-supervisor-contact.png}
        \caption{มุมมอง Supervisor}
    \end{subfigure}
    \begin{subfigure}[b]{0.49\textwidth}
        \centering
        \includegraphics[width=\linewidth]{chapter3/dashboard-hmi/role-observer-contact.png}
        \caption{มุมมอง Observer}
    \end{subfigure}
    \caption{ภาพหน้าจอระบบในมุมมอง Supervisor และ Observer สำหรับการติดตามงานและสถานะตามสิทธิ์}
    \label{fig:ch3_hmi_supervisor_observer_contact}
\end{figure}
\FloatBarrier

หน้าจอ Patient ถูกออกแบบให้เห็นเฉพาะข้อมูลของผู้ใช้รายนั้น เช่น care plan, schedule, medication, room control, service request และช่องทางติดต่อผู้ดูแล ไม่แสดงรายการผู้พักอาศัยคนอื่นหรือข้อมูลบุคลากรที่ไม่เกี่ยวข้อง ภาพที่~\ref{fig:ch3_hmi_patient_dashboard} เป็นภาพหน้าจอของ Patient dashboard ที่ใช้ในการตรวจสอบรอบปัจจุบัน ส่วนภาพที่~\ref{fig:ch3_hmi_patient_contact} แสดงหน้าจออื่นใน portal ของผู้ป่วย เช่น schedule, pharmacy, room controls, services, messages, support และ settings

\begin{figure}[!htbp]
    \centering
    \includegraphics[width=0.9\textwidth]{chapter3/dashboard-hmi/patient-dashboard-overview.png}
    \caption{ภาพหน้าจอ Patient dashboard สำหรับผู้ใช้รายบุคคล แสดงข้อมูล care plan สถานะรถเข็น และช่องทางขอความช่วยเหลือ}
    \label{fig:ch3_hmi_patient_dashboard}
\end{figure}
\FloatBarrier

\begin{figure}[!htbp]
    \centering
    \includegraphics[width=0.98\textwidth]{chapter3/dashboard-hmi/role-patient-contact.png}
    \caption{ภาพหน้าจอระบบในมุมมอง Patient ครอบคลุม personal dashboard, schedule, pharmacy, room controls, services, messages, support และ settings}
    \label{fig:ch3_hmi_patient_contact}
\end{figure}
\FloatBarrier

EaseAI เป็น assistant ที่เปิดจาก dashboard เพื่อช่วยผู้ดูแลค้นข้อมูล สรุปบริบท หรือเตรียม action candidate จากข้อมูลที่ระบบอนุญาตให้เข้าถึง การใช้งานเริ่มจากผู้ใช้เปิดแผง EaseAI เลือกคำถามตัวอย่างหรือพิมพ์คำถามเอง จากนั้นระบบส่งข้อความพร้อม role, workspace และบริบทที่จำเป็นเข้าสู่ AI pipeline ในหัวข้อ~\ref{sec:ch3_ai} หากคำตอบเป็นเพียงการสรุปข้อมูล ระบบจะแสดงผลกลับในหน้าต่างสนทนา แต่ถ้าเกี่ยวข้องกับการสั่งงานจริงจะต้องผ่าน execution gate และเงื่อนไขการยืนยันก่อนเสมอ ภาพที่~\ref{fig:ch3_hmi_easeai} แสดงตัวอย่างหน้าจอ EaseAI บนมุมมอง Admin

\begin{figure}[!htbp]
    \centering
    \includegraphics[width=0.9\textwidth]{chapter3/dashboard-hmi/admin-easeai-chat.png}
    \caption{ภาพหน้าจอ EaseAI ใน dashboard สำหรับถามข้อมูลหรือเตรียม action ภายใต้ role permission และ MCP execution gate}
    \label{fig:ch3_hmi_easeai}
\end{figure}
\FloatBarrier

\section{ระบบ AI และ MCP Tool Pipeline}
\label{sec:ch3_ai}
ส่วน AI ของระบบถูกออกแบบใหม่จากปัญหาของแนวทางเดิมที่ให้ LLM และ MCP agent จัดการงานส่วนใหญ่ในเส้นทางเดียว เมื่อระบบมี tool หลายกลุ่ม เช่น การค้นข้อมูลผู้พักอาศัย การอ่านข้อมูล sensor การสรุปสถานะ และการเตรียม action ให้ผู้ดูแลตรวจสอบ prompt ต้องบรรจุคำอธิบาย tool, schema, policy และบริบทจำนวนมาก ส่งผลให้เกิด token overhead สูงและ latency เพิ่มขึ้นมาก จากบันทึกการพัฒนาระบบเดิมพบว่าโมเดลขนาด 7B บางรอบใช้เวลาประมาณ 30--45 นาที หรือมากกว่านั้นเมื่อต้อง reasoning กับ tool หลายขั้นตอน แม้ภายหลังจะลดขนาดโมเดลและปรับ prompt แล้ว เวลาตอบสนองยังไม่เหมาะกับ dashboard ที่ผู้ดูแลต้องรอผลในรอบการทำงานสั้นกว่าเดิม

หลังการทดสอบรอบแรก ระบบถูกปรับเป็นแนวทางที่แยก logic เชิงสัญลักษณ์ออกจากการสังเคราะห์ภาษาของ LLM โดยให้ deterministic rule, schema, policy และ tool contract ทำงานร่วมกับ LLM แทนการปล่อยให้โมเดลตัดสินใจทุกขั้นตอน แนวทางนี้เป็นการออกแบบแบบ neuro-symbolic-inspired กล่าวคือใช้ LLM สำหรับความเข้าใจเชิงภาษาและการเรียบเรียงคำตอบ ส่วนเงื่อนไขเชิงระบบและการสั่งงานจริงยังถูกควบคุมด้วยกฎและ schema ที่ตรวจสอบได้ จึงไม่กล่าวอ้างว่าเป็นระบบ neuro-symbolic เต็มรูปแบบ แต่ใช้หลักการดังกล่าวเพื่อลด hallucination และลดคำสั่งที่อยู่นอกขอบเขต

pipeline หลักของระบบแบ่งเป็น 5 layer ได้แก่ Layer 1 Intent Router, Layer 2 Context Assembler, Layer 3 Behavioral State, Layer 4 Constrained LLM และ Layer 5 Execution Gate อย่างไรก็ตาม Layer 3 ไม่ได้อยู่ใน execute path หลักแบบ synchronous เสมอไป เส้นทางหลักของคำสั่งคือ Layer 1 \(\rightarrow\) Layer 2 \(\rightarrow\) Layer 4 \(\rightarrow\) Layer 5 ส่วน Layer 3 ทำหน้าที่เตรียมและเก็บ behavioral context จากข้อมูลผู้ใช้หรือเหตุการณ์ เพื่อให้ระบบนำบริบทนี้กลับมาใช้ในคำถามหรือการประเมินรอบถัดไปได้

\begin{figure}[!htbp]
    \centering
    \resizebox{0.92\textwidth}{!}{%
    \begin{tikzpicture}[
        box/.style={draw, rounded corners, align=center, text width=2.55cm, minimum height=0.88cm, font=\scriptsize, fill=white},
        side/.style={draw, rounded corners, align=center, text width=2.8cm, minimum height=0.88cm, font=\scriptsize, fill=yellow!12},
        arrow/.style={->, >=Stealth, thick}
    ]
        \node[box, fill=green!10] (l1) {Layer 1\\Intent Router};
        \node[box, fill=green!10, right=0.48cm of l1] (l2) {Layer 2\\Context Assembler};
        \node[box, fill=green!10, right=0.48cm of l2] (l4) {Layer 4\\Constrained LLM};
        \node[box, fill=green!10, right=0.48cm of l4] (l5) {Layer 5\\Execution Gate};
        \node[side, below=0.72cm of l2] (l3) {Layer 3\\Behavioral State};
        \draw[arrow] (l1) -- (l2);
        \draw[arrow] (l2) -- (l4);
        \draw[arrow] (l4) -- node[above, font=\tiny] {candidate} (l5);
        \draw[arrow, dashed] (l2.south) -- node[right, font=\tiny] {event} (l3.north);
        \draw[arrow, dashed] (l3.west) -- ++(-0.55,0) |- node[pos=0.28, left, font=\tiny] {context} (l1.south);
    \end{tikzpicture}}
    \caption{โครงสร้าง AI pipeline 5 layer โดย Layer 3 ทำหน้าที่เตรียม behavioral context นอกเส้นทาง execute หลัก}
    \label{fig:ch3-ai-flow}
\end{figure}
\FloatBarrier

\textbf{Layer 1 Intent Router} ทำหน้าที่รับคำสั่งจากผู้ใช้แล้วจัดกลุ่มเจตนาเบื้องต้นก่อนให้ LLM ทำงาน ระบบเริ่มจาก deterministic matching เช่น keyword, route, role และ schema ที่คาดหวัง หาก confidence เพียงพอจะส่ง query พร้อม intent ไปยัง Layer 2 ได้ทันที หากยังไม่ชัดเจนจึงใช้ semantic matching เพื่อวัดความใกล้เคียงเชิงความหมาย และหากยังขาดข้อมูลสำคัญ เช่น ไม่ระบุผู้พักอาศัย ห้อง หรือ action ที่ต้องการ ระบบจะหยุด pipeline แล้วถามกลับแทนการเดาคำสั่ง

\begin{figure}[!htbp]
    \centering
    \resizebox{0.92\textwidth}{!}{%
    \begin{tikzpicture}[
        box/.style={draw, align=center, text width=3.1cm, minimum height=0.78cm, font=\scriptsize, fill=white},
        arrow/.style={->, >=Stealth, thick}
    ]
        \node[font=\scriptsize] (q) {Query};
        \node[box, right=0.45cm of q] (det) {\textbf{Deterministic Matching}\\keyword, route, schema};
        \node[box, right=0.58cm of det] (sem) {\textbf{Semantic Matching}\\embedding similarity};
        \node[box, right=0.58cm of sem] (clar) {\textbf{Deterministic Clarification}\\ask missing facts};
        \node[font=\scriptsize, below=0.65cm of sem] (out) {Send query + intent to Layer 2};
        \draw[arrow] (q) -- (det);
        \draw[arrow] (det) -- node[above, font=\tiny] {else} (sem);
        \draw[arrow] (sem) -- node[above, font=\tiny] {else} (clar);
        \draw[arrow] (det.south) |- node[pos=0.23, left, font=\tiny] {if confident} (out.west);
        \draw[arrow] (sem.south) -- (out.north);
        \draw[arrow] (clar.south) -- ++(0,-0.35) node[below, font=\tiny, align=center] {break pipeline\\and ask clarification};
    \end{tikzpicture}}
    \caption{Layer 1 Intent Router สำหรับจับเจตนาแบบ deterministic ก่อนใช้ semantic matching และ clarification}
    \label{fig:ch3_ai_intent_router}
\end{figure}
\FloatBarrier

\textbf{Layer 2 Context Assembler} รับ intent จาก Layer 1 แล้วเลือกบริบทที่เกี่ยวข้องกับคำสั่ง เช่น ข้อมูลผู้พักอาศัย สถานะ sensor ล่าสุด ประวัติ telemetry สรุปเชิงพฤติกรรม หรือเอกสารใน RAG ขั้นตอนนี้ไม่ได้ดึงข้อมูลทั้งหมดเข้า prompt แต่เริ่มจากการวางแผนว่าคำสั่งต้องใช้หัวข้อใด จากนั้น fetch และ rank เฉพาะ chunk ที่อนุญาตตาม role และ policy ก่อนแพ็กบริบทให้สั้น ไม่ซ้ำ และตรวจย้อนกลับได้

\begin{figure}[!htbp]
    \centering
    \resizebox{0.94\textwidth}{!}{%
    \begin{tikzpicture}[
        box/.style={draw, align=center, text width=2.8cm, minimum height=0.8cm, font=\scriptsize, fill=white},
        arrow/.style={->, >=Stealth, thick}
    ]
        \node[font=\scriptsize, align=center] (intent) {Query\\intent};
        \node[box, right=0.45cm of intent] (plan) {\textbf{Plan}\\choose topic area};
        \node[box, right=0.45cm of plan] (rank) {\textbf{Fetch + Rank}\\allowed chunks};
        \node[box, right=0.45cm of rank] (pack) {\textbf{Pack}\\small clean context};
        \node[box, right=0.45cm of pack] (gate) {\textbf{Gate + Context}\\schema ready};
        \node[font=\scriptsize, align=center, right=0.45cm of gate] (ctx) {Query\\context};
        \draw[arrow] (intent) -- (plan);
        \draw[arrow] (plan) -- (rank);
        \draw[arrow] (rank) -- (pack);
        \draw[arrow] (pack) -- (gate);
        \draw[arrow] (gate) -- (ctx);
    \end{tikzpicture}}
    \caption{Layer 2 Context Assembler สำหรับดึงและจัดแพ็กบริบทก่อนส่งเข้า LLM}
    \label{fig:ch3_ai_context_assembler}
\end{figure}
\FloatBarrier

บริบทที่ถูกส่งต่อจะประกอบด้วยข้อมูลที่จำเป็นต่อคำตอบหรือการสั่งงานเท่านั้น เช่น patient id, device id, ค่า telemetry ล่าสุด, สิทธิ์ของผู้ใช้ และข้อความอ้างอิงที่ถูกดึงจาก retrieval logic วิธีนี้ทำให้ Layer 4 ไม่ต้องรับ prompt ขนาดใหญ่ทุกครั้ง และลดโอกาสที่ LLM จะเห็นข้อมูลที่ไม่เกี่ยวข้องจนสรุปผิด

\begin{figure}[!htbp]
    \centering
    \resizebox{0.88\textwidth}{!}{%
    \begin{tikzpicture}[
        box/.style={draw, align=center, text width=3.0cm, minimum height=0.8cm, font=\scriptsize, fill=white},
        arrow/.style={->, >=Stealth, thick}
    ]
        \node[font=\scriptsize, align=center] (ctx) {Query\\context};
        \node[box, right=0.45cm of ctx] (schema) {\textbf{Schema-Bound Prompting}\\tool schema + prompt};
        \node[box, right=0.55cm of schema] (mcp) {\textbf{MCP Client Inference}\\LLM via MCP client};
        \node[box, right=0.55cm of mcp] (handoff) {\textbf{Execution Candidate Handoff}\\package for Layer 5};
        \node[font=\scriptsize, align=center, right=0.45cm of handoff] (cand) {Execution\\candidate};
        \draw[arrow] (ctx) -- (schema);
        \draw[arrow] (schema) -- (mcp);
        \draw[arrow] (mcp) -- (handoff);
        \draw[arrow] (handoff) -- (cand);
    \end{tikzpicture}}
    \caption{Layer 4 Constrained LLM สำหรับสังเคราะห์คำตอบหรือ execution candidate ภายใต้ schema ที่กำหนด}
    \label{fig:ch3_ai_constrained_llm}
\end{figure}
\FloatBarrier

\textbf{Layer 4 Constrained LLM} เป็นชั้นที่ใช้ LLM สังเคราะห์คำตอบภายใต้ prompt และ schema ที่จำกัดไว้ ระบบใช้ schema-bound prompting เพื่อระบุรูปแบบผลลัพธ์ที่ต้องการ เช่น ต้องตอบเป็นข้อความสรุป ต้องคืนค่า action candidate หรือไม่ควรเรียก tool ใดในสถานการณ์นั้น หากคำสั่งเกี่ยวข้องกับ MCP ระบบจะให้ LLM สร้าง execution candidate ที่มี argument ตาม schema เท่านั้น ยังไม่ให้ดำเนินการกับระบบจริงในชั้นนี้

\begin{figure}[!htbp]
    \centering
    \resizebox{0.86\textwidth}{!}{%
    \begin{tikzpicture}[
        box/.style={draw, align=center, text width=3.0cm, minimum height=0.8cm, font=\scriptsize, fill=white},
        arrow/.style={->, >=Stealth, thick}
    ]
        \node[font=\scriptsize, align=center] (cand) {tool call\\args package};
        \node[box, right=0.55cm of cand] (val) {\textbf{Validate and Authorize}\\schema + policy check};
        \node[box, right=0.75cm of val] (exe) {\textbf{Execute}\\send command to MCP server};
        \node[font=\scriptsize, align=center, right=0.55cm of exe] (res) {execution\\user-visible result};
        \draw[arrow] (cand) -- (val);
        \draw[arrow] (val) -- (exe);
        \draw[arrow] (exe) -- (res);
    \end{tikzpicture}}
    \caption{Layer 5 Execution Gate สำหรับตรวจ schema, policy และยืนยันการสั่งงานผ่าน MCP tool}
    \label{fig:ch3_ai_execution_gate}
\end{figure}
\FloatBarrier

\textbf{Layer 5 Execution Gate} เป็นชั้นป้องกันก่อนส่งคำสั่งไปยัง MCP server หรือ backend tool โดยตรวจ schema, policy, role, allowlist และเงื่อนไขด้านความปลอดภัยของ action อีกครั้ง หาก argument ไม่ครบ สิทธิ์ไม่ตรง หรือ action เสี่ยงเกินขอบเขต ระบบจะไม่ execute และส่งผลกลับเป็นคำถามหรือคำอธิบายให้ผู้ดูแลตรวจสอบ วิธีนี้ทำให้ LLM ไม่สามารถสั่งงานจริงได้โดยลำพัง แม้จะสร้างข้อความที่ดูเหมือนคำสั่งถูกต้อง

\textbf{Layer 3 Behavioral State} ทำงานเป็นกระบวนการเตรียมบริบทจากข้อมูล biometrics, movement, resident id และ timestamp ที่ระบบรับมาอย่างต่อเนื่อง ข้อมูลรายวันถูกบีบให้เป็นสัญญาณที่อ่านได้ เช่น ค่าเฉลี่ย ส่วนเบี่ยงเบนมาตรฐาน rolling window แนวโน้มจาก linear regression และจำนวนเหตุการณ์ผิดปกติ จากนั้นระบบใช้ rule หรือ model เช่น XGBoost ในระดับต้นแบบเพื่อจัดกลุ่มความเร่งด่วนและแกนปัญหาที่ควรเฝ้าดู ผลลัพธ์ถูกเก็บเป็น behavioral database เพื่อให้ Layer 2 ดึงกลับไปประกอบคำตอบรอบถัดไป ไม่ใช่การวินิจฉัยทางการแพทย์หรือการตัดสินแทนผู้ดูแล

\begin{figure}[!htbp]
    \centering
    \resizebox{0.92\textwidth}{!}{%
    \begin{tikzpicture}[
        box/.style={draw, align=center, text width=2.95cm, minimum height=0.8cm, font=\scriptsize, fill=white},
        arrow/.style={->, >=Stealth, thick}
    ]
        \node[font=\scriptsize, align=center] (in) {biometrics\\movement\\resident id\\timestamp};
        \node[box, right=0.5cm of in] (stat) {\textbf{Statistical Analysis}\\mean, STD, trend};
        \node[box, right=0.55cm of stat] (risk) {\textbf{Risk Analysis}\\rule + model signals};
        \node[box, right=0.55cm of risk] (store) {\textbf{Storing}\\queryable history};
        \node[font=\scriptsize, align=center, right=0.5cm of store] (out) {behavioral DB\\updates};
        \draw[arrow] (in) -- (stat);
        \draw[arrow] (stat) -- (risk);
        \draw[arrow] (risk) -- (store);
        \draw[arrow] (store) -- (out);
    \end{tikzpicture}}
    \caption{Layer 3 Behavioral State สำหรับเตรียมบริบทพฤติกรรมจากข้อมูล sensor และเหตุการณ์ของผู้ใช้}
    \label{fig:ch3_ai_behavioral_state}
\end{figure}
\FloatBarrier

ส่วน Edge AI ที่เกี่ยวข้องกับ Layer 3 ยังอยู่ในระดับ proof of concept โดยแนวคิดด้าน motion anomaly ใช้ข้อมูล IMU เช่น gyro magnitude, acceleration magnitude และ rolling standard deviation เพื่อสังเกตการเคลื่อนไหวที่ต่างจากรูปแบบปกติ ส่วนแนวคิดด้าน health anomaly ใช้ HR, RR interval และ PPG จาก Polar Verity Sense เป็นบริบทประกอบ หากพบเหตุการณ์ที่เข้าข่ายผิดปกติ ระบบจะสร้าง event หรือ context ให้ผู้ดูแลตรวจสอบผ่าน dashboard หรือใช้ประกอบการตอบของ AI ต่อไป ไม่ได้ใช้เป็นผลตัดสินทางคลินิก

\begin{figure}[!htbp]
    \centering
    \resizebox{0.88\textwidth}{!}{%
    \begin{tikzpicture}[
        box/.style={draw, align=center, text width=2.8cm, minimum height=0.78cm, font=\scriptsize, fill=white},
        arrow/.style={->, >=Stealth, thick}
    ]
        \node[font=\scriptsize, align=center] (in) {I/O turn\\payload};
        \node[box, right=0.5cm of in] (tag) {\textbf{Tagging}\\session + I/O};
        \node[box, right=0.55cm of tag] (embed) {\textbf{Embedding Signals}\\rank + newness};
        \node[box, right=0.55cm of embed] (policy) {\textbf{Policy}\\limit active sessions};
        \node[font=\scriptsize, align=center, right=0.5cm of policy] (out) {context window\\update};
        \draw[arrow] (in) -- (tag);
        \draw[arrow] (tag) -- (embed);
        \draw[arrow] (embed) -- (policy);
        \draw[arrow] (policy) -- (out);
    \end{tikzpicture}}
    \caption{Context Window Management สำหรับจำกัดบริบทให้มีขนาดเหมาะสมและตรวจสอบย้อนกลับได้}
    \label{fig:ch3_ai_context_window}
\end{figure}
\FloatBarrier

Context Window Management เป็นส่วนสนับสนุนที่ช่วยให้การสนทนาและผลลัพธ์จาก tool ไม่โตจนเกินขอบเขตของโมเดล แต่ยังย้อนกลับได้ว่าแต่ละคำตอบใช้ข้อมูลใด ระบบจึง tag payload ตาม session และ I/O turn เก็บ embedding signal สำหรับใช้ rank ร่วมกับความใหม่ของข้อมูล และมี policy จำกัดจำนวน active session หรือ bundle ที่เก็บใน context เมื่อข้อมูลเก่าเกินจำเป็นจะถูกลดความสำคัญหรือถอดออกจาก prompt แต่ยังคงบันทึกในฐานข้อมูลหรือ log เพื่อการตรวจสอบภายหลัง

MCP ทำหน้าที่เป็นสัญญาระหว่าง AI runtime กับระบบจริง โดยฝั่ง agent runtime เป็นผู้ประกอบเจตนา บริบท และ schema ก่อนส่ง execution candidate ไปยัง execution gate ส่วน MCP server เปิด resources, prompts และ tools ที่ผ่านการควบคุมสิทธิ์แล้ว การแยกเช่นนี้ช่วยลดโอกาสที่ LLM จะเรียกใช้คำสั่งนอกขอบเขต เพราะ tool ที่รันจริงต้องผ่านทั้ง schema validation, allowlist ตามบทบาท และ execution gate ก่อนเสมอ

\section{แนวทางการประเมินผลระบบ}
\label{sec:ch3_evaluation_plan}
การประเมินผลระบบถูกออกแบบให้เชื่อมกับองค์ประกอบหลักของสถาปัตยกรรมที่อธิบายในบทนี้ โดยไม่สรุปผลเกินกว่าข้อมูลทดลองที่มี ในบทที่ 4 แต่ละหัวข้อจึงระบุสภาพแวดล้อมการทดสอบ วิธีทดสอบ ผลที่ได้ และการวิเคราะห์ในรูปแบบย่อหน้าต่อเนื่อง เพื่อให้ผู้อ่านแยกได้ว่าปัญหาหรือข้อจำกัดเกิดจาก sensor, gateway, server, dashboard หรือ AI pipeline

\begin{table}[H]
    \centering
    \caption{แนวทางการประเมินผลที่ใช้ต่อเนื่องไปยังบทผลการทดลอง}
    \label{tab:ch3_evaluation_plan}
    \small
    \renewcommand{\arraystretch}{1.22}
    \begin{tabular}{|L{3.4cm}|L{4.8cm}|L{4.5cm}|}
        \hline
        \textbf{ส่วนที่ประเมิน} & \textbf{สิ่งที่ทดสอบ} & \textbf{ตัวชี้วัดหรือผลที่ต้องรายงาน} \\ \hline
        Sensor และอุปกรณ์ภาคสนาม & M5StickC Plus2 IMU, battery, Polar HR/RR/PPG, BLE connection & ความต่อเนื่องของข้อมูล ความถูกต้องเชิงสังเกต และความเสถียรของการเชื่อมต่อ \\ \hline
        Indoor localization & การเก็บ fingerprint, KNN, max-RSSI fallback และการส่งผล prediction & accuracy, error หรือความคลาดเคลื่อนรายจุด พร้อมวิเคราะห์ RSSI fluctuation \\ \hline
        Mobile gateway และ server & BLE pairing, MQTT/REST publish, reconnect/offline handling, backend ingest & อัตราการส่งข้อมูลสำเร็จ latency และการแสดงผลบน dashboard \\ \hline
        Dashboard/HMI & workflow ของ Admin, Head Nurse, Supervisor, Observer และ Patient & task completion, ข้อสังเกตจากการใช้งาน และ SUS Score สำหรับ usability \\ \hline
        AI/MCP pipeline & Intent routing, context retrieval, constrained LLM, execution gate และ trace ต่อ layer & IRR, cosine similarity, latency ต่อ layer และความปลอดภัยของ tool execution \\ \hline
    \end{tabular}
\end{table}

การเพิ่ม SUS Score ในการประเมินช่วยให้ส่วน HMI ไม่ได้ประเมินจากภาพหน้าจอหรือความเห็นทั่วไปเพียงอย่างเดียว แต่มีคะแนนมาตรฐานด้าน usability ประกอบกับ task scenario ของผู้ใช้ ขณะที่การประเมิน AI ใช้ทั้งความเห็นของมนุษย์ต่อคำตอบ ความใกล้เคียงเชิงข้อความ และ latency ราย layer เพื่อแสดงให้เห็นว่าการแยก logic ออกจาก LLM ช่วยให้ระบบตรวจสอบและปรับปรุงได้เป็นส่วน ๆ
